% ============================================================
%  chapters/04_discussion.tex
%  Evidence base: 78 included experimental studies
% ============================================================

\section{Discussion}
\label{sec:discussion}
\label{sec:discussion:findings}

\subsection{Principal Findings}

This systematic review of 78 histopathology AI fairness studies identifies three converging evidence patterns. First, pathology foundation models faithfully encode TSS signatures despite claims of domain invariance~\citep{vorontsov2024foundation, campanella2025clinical}. Linear probes recover TSS at $>$90\% accuracy from frozen foundation-model features~\citep{komen2024batch}, corresponding to site-classification AUROC $>$0.96 on TCGA~\citep{howard2021site}. These complementary metrics (accuracy and AUROC) both confirm strong TSS encoding in learned features. When TSS correlates with labels, downstream classification collapses~\citep{howard2021site, dehkharghanian2023biased}. Feature-level corrections such as ComBat~\citep{murchan2024combat} offer partial relief, but the persistence of TSS encoding underscores the need for fairness-aware pretraining rather than post-hoc correction. Cross-modal evidence from radiology confirms this is a structural property of medical imaging pipelines, not pathology-specific~\citep{gao2022race, yang2025demographic}.

Second, fairness metrics are systematically under-reported: only 19.2\% (95\% CI: 12.0--29.3\%) of studies report any fairness metric and 7.7\% (95\% CI: 3.6--16.0\%) report worst-group accuracy. This reporting asymmetry masks clinically significant disparities: overall AUROC of 0.98 coexists with 3--16\% subgroup gaps~\citep{vaidya2024demographic}. Third, the evidence base is structurally concentrated on TCGA (42.3\%, 95\% CI: 32.0--53.2\% of studies), which embeds strong site--label dependencies and demographic skews~\citep{howard2021site, wang2018tcga}. These three patterns reinforce one another: site leakage goes undetected without fairness metrics, and both deficits are amplified by TCGA's dominance.

Despite this consensus, every normative recommendation exceeds realised practice by 55--89 percentage points (Table~\ref{tab:consensus_vs_practice}). Four factors sustain this gap: recommendations are costless to reviewers but costly to implementers; compliant datasets largely do not exist~\citep{norori2021addressing, ceccon2025underrepresentation}; reviewer consensus is trailing-edge (29.5\% of experiments pre-date 2023); and the reward structure privileges aggregate AUROC over subgroup audits~\citep{cheng2021challenges, chen2023algorithmic, yang2024survey}.

%% ─── Consensus vs Practice table ─────────────────────────
\begin{table}[htbp]
  \centering
  \caption{Review consensus ($\geq 80\%$) versus realised practice.
           Gaps in percentage points.}
  \label{tab:consensus_vs_practice}
  \small
  \begin{tabularx}{\linewidth}{@{}X c c c@{}}
    \toprule
    \textbf{Recommendation} & \textbf{Consensus} &
      \textbf{Practice} & \textbf{Gap} \\
    \midrule
    Multi-site training data~\citep{riccilara2022addressing, montezuma2025unbiased, howard2021site}
      & $\sim$95\% & 18.7\% & $\sim$76 pp \\
    External validation~\citep{montezuma2025unbiased, chen2023algorithmic, lin2025contrastive}
      & $\sim$95\% & 29.5\% & $\sim$66 pp \\
    Demographic metadata~\citep{chen2023algorithmic, sounderajah2021stardai, vaidya2024demographic}
      & $\sim$85\% & $\leq$28.2\% & $\sim$57 pp \\
    Batch-effect mitigation~\citep{komen2024batch, chinta2025aidriven, murchan2024combat}
      & $\sim$95\% & $\sim$6\%\textsuperscript{*} & $\sim$89 pp \\
    Standardised fairness reporting~\citep{sounderajah2021stardai, chen2023algorithmic, zong2023medfair}
      & $\sim$82\% & 19.2\% & $\sim$63 pp \\
    Preserved-site CV~\citep{howard2021site, chinta2025aidriven, kheiri2025investigation}
      & $\sim$90\% & $\sim$35\% & $\sim$55 pp \\
    Worst-group accuracy/AUC~\citep{zong2023medfair, chen2023algorithmic, soltan2024challenges}
      & $\sim$65\% & 7.7\% & $\sim$57 pp \\
    \bottomrule
  \end{tabularx}
  \textsuperscript{*}Feature-level methods only (ComBat, FAIR-Path, FLEX, etc.); approximately 5 studies; stain normalization excluded (fails to mitigate).
\end{table}

\subsection{Key Challenges and Evidence Gaps}
\label{sec:discussion:patterns}

\subsubsection{Measurement Asymmetry and Bias Taxonomy}
Research attention does not track clinical consequence, which represents a critical misalignment in the current literature. Technical and domain-shift bias, along with selection and sampling bias, each attracted eight studies, the highest counts in our corpus, while demographic and algorithmic bias trailed at five studies each~(Figure~\ref{fig:bias_taxonomy}; Table~\ref{tab:bias_taxonomy})~\citep{vaidya2024demographic, howard2021site, chen2025predicting, lin2025contrastive, komen2024batch, dehkharghanian2023biased, kheiri2025investigation, ling2024agent, campanella2019clinical, zong2023medfair, cecconi2024fairness, soltan2024challenges, celi2022sources, musa2026cross, tasci2022bias, ktena2024generative, hagele2020resolving, vanbooven2025mitigating}. Cross-modal evidence from radiology and dermatology corroborates histopathology-specific findings throughout this section, confirming these are structural properties of medical imaging pipelines.

This asymmetry reflects the relative ease of measuring technical proxies versus the infrastructural requirements for demographic fairness evaluation. Yet demographic biases produce the largest measured disparities: studies report 3--16\% AUROC gaps across 3~cancer diagnostic tasks~\citep{vaidya2024demographic}, though cross-population mutation-rate variation provides a partial biological basis for these differences~\citep{polin2025mutation}. Other measured disparities include up to 33-percentage-point accuracy differences across racial groups in cohorts with $>$70\% White patients~\citep{soltan2024challenges}, 5--6\% lower Dice scores for minority ethnic groups in cardiac MRI~\citep{puyolantion2022fairness}, and algorithmic bias in NLP-based biomarker classification~\citep{tschida2025evaluating}. These disparities substantially exceed the $F_1$ variation of 0.09--0.27 across scanners~\citep{aubreville2022mitosis} and the cross-batch AUROC degradation of 0.22--0.36~\citep{lin2025stain} typically reported for technical biases.

This asymmetry stems from measurement infrastructure: technical bias can be quantified from site-classification accuracy and cross-batch AUROC without demographic metadata~\citep{dehkharghanian2023biased, kheiri2025investigation}, whereas demographic fairness requires stratified evaluation against metadata that most datasets simply do not provide. The field studies what it can measure, not what matters most for patients.

This measurement asymmetry is not unique to histopathology. Cross-domain evidence from radiology reveals similar patterns: a radiology-focused review identified 29 bias sources across five pipeline stages (data collection, annotation, model development, evaluation, and deployment)~\citep{drukker2023toward}, yet its experimental operationalisation similarly reduces to a handful of measurable proxies.

We see a consistent divergence between review taxonomies, which are broad and top-down~\citep{yang2024survey, chen2023algorithmic, norori2021addressing, drukker2023toward}, and experimental operationalisation, which relies on narrow proxies such as site accuracy, mutual information, subgroup AUROC gap, and shortcut-attribution scores~\citep{ly2024shortcut, dawood2026confounding, banerjee2023shortcuts}. This divergence carries three consequences. First, normative consensus runs ahead of the empirical base. Second, review categories collapse mechanistically distinct biases that actually require different mitigations. Third, the hardest-to-measure categories are also the hardest to mitigate~\citep{zong2023medfair, cecconi2024fairness}.

TCGA demographics illustrate this tension: the dataset's racial composition skews toward non-Hispanic Whites, with Black patients under-represented in 8 of 13 audited cancer types~\citep{wang2018tcga}, and its age distribution runs 3.9 years younger than the US general cancer population on average~\citep{wang2018tcga}.

Systematic underdiagnosis represents a further evidence gap, compounded by regulatory failures. Supervised chest radiograph classifiers underdiagnose Black, Hispanic, female, and Medicaid-insured patients, with intersectional subgroups showing underdiagnosis rates exceeding 50\%~\citep{seyyedkalantari2021underdiagnosis}. Demographic shortcuts compound these disparities: models encoding such shortcuts show a correlation of $R{=}0.82$ between demographic-encoding strength and subgroup unfairness. Locally optimal corrections fail under real-world distribution shift, producing false-negative-rate gaps of 30\% across age groups~\citep{yang2024limits}. Even after algorithmic correction, models may retain encoded demographic biases. Fairness-aware methods may not generalise across datasets or tasks~\citep{yang2024limits}. Pre-market regulatory oversight has similarly failed to enforce transparency: a scoping review of 692 FDA-approved AI/ML medical devices from 1995--2023 found that only 3.6\% reported race or ethnicity and 99.1\% provided no socioeconomic data, establishing that pre-market authorisation under existing frameworks has not enforced demographic transparency~\citep{muralidharan2024fda}. Privacy-preserving federated approaches that maintain fairness present a further challenge. Fairness-aware aggregation beyond Prop-FFL~\citep{hosseini2023proportionally} remains underdeveloped, encompassing federated fairness auditing without centralising data and privacy-preserving substrates calibrated to avoid erasing subgroup disparities~\citep{lu2022federated}, though FlexFair addresses this directly by jointly optimising demographic parity, equalised odds, and accuracy in federated medical imaging without centralising data~\citep{xing2025flexfair}.

\subsubsection{Mitigation Effectiveness}
\label{sec:discussion:effectiveness}
Image-level stain normalization remains the dominant experimental approach, yet conventional methods consistently fail to remove TSS-level confounding. Cross-batch AUROC still collapses after Reinhard, Macenko, and Vahadane correction~\citep{komen2024batch}. TSS-prediction AUROC remains above 0.85 after normalization~\citep{howard2021site}, confirming the TSS encoding patterns documented in Section~\ref{sec:discussion:findings}. Center identity can still be recovered at 50\% accuracy after greyscale normalization, down only modestly from 59\% uncorrected~\citep{kheiri2025investigation}. Post-2023 reviews have recalibrated, with the majority of recent systematic reviews now characterizing conventional stain normalization alone as insufficient~\citep{montezuma2025unbiased, chinta2025aidriven}, yet it remains the modal experimental mitigation. Feature-level and fairness-aware methods prove more effective but receive less endorsement. ComBat reduced TSS-prediction AUROC to chance while preserving MSI prediction and improving external BRAF detection~\citep{murchan2024combat}. FAIR-Path mitigated 88.5--91.1\% of disparities across 20 TCGA cancer types (88.5\% internal, 91.1\% external validation; range reflects validation setting)~\citep{lin2025contrastive}, FLEX reduced the fairness gap by $\sim$61\% while improving OOD AUROC by 6.4--7.2\%~\citep{huang2025knowledge}, and SPARE improved minority-group accuracy by 3.7--5.8\%~\citep{xu2026spare}. Yet combined fairness-aware endorsement reaches only $\sim$65\%, falling below the 80\% threshold typical for reporting standards.

Cross-modal evidence from radiology sets a sobering limit. Even after algorithmic correction, models encoding demographic shortcuts show a correlation of $R{=}0.82$ between demographic-encoding strength and subgroup unfairness, and locally optimal corrections fail under real-world distribution shift, producing false-negative-rate gaps of 30\% across age groups~\citep{yang2024limits}. Natural-image domain-generalisation baselines fare even worse: CORAL, IRM, and Group DRO leave a persistent 22.9-point gap on WILDS-CAMELYON17~\citep{koh2021wilds}, while histopathology-tuned domain-adaptation methods such as AIDA have eclipsed them~\citep{asadiaghbolagh2024learning}. Foundation models present a similar dual character. They improve representation quality~\citep{ciga2023self, huang2025knowledge} but faithfully encode every systematic signal in training data, including TSS signatures~\citep{komen2024batch, dehkharghanian2023biased, sikaroudi2023generalization}. Their fairness properties are determined not by the pretraining procedure itself, but by what happens after pretraining: feature-level harmonisation, fairness-aware fine-tuning, and preserved-site evaluation. This post-hoc dependence highlights the need for fairness-aware pretraining objectives rather than reliance on downstream correction.

The most widely endorsed practices, namely preserved-site cross-validation, external validation, and demographic reporting (each at $\sim$85--90\% endorsement), constitute audit infrastructure rather than bias reduction. The field is being directed to instrument bias more reliably, not to remove it. This creates a paradox: we have built sophisticated tools to measure our failures (Section~\ref{sec:discussion:findings}) but lack the incentives and infrastructure to prevent them.

\subsubsection{Data Infrastructure Deficits}
\label{sec:discussion:data}
A 39-dataset audit reveals stark asymmetry in data infrastructure that shapes what the field can study. Technical metadata (modality, staining, cancer type) is universally available (97\% or higher), but race/ethnicity, sex, and age are documented in only 28.2\%, 25.6\%, and 23.1\% of datasets, respectively. Six fairness-critical fields have zero availability across all audited datasets: genetic ancestry, Fitzpatrick skin tone, intersectional identifiers, preferred language, protected-attribute flags, and data-quality scores~\citep{norori2021addressing, celi2022sources, vaidya2024demographic}. This matters because dataset composition, not model design, determines classifier fairness: training on sex-imbalanced medical imaging datasets produces significantly lower performance for underrepresented genders even without any explicit bias mechanism~\citep{larrazabal2020gender}. Demographic metadata concentrates precisely where batch effects are strongest. Race appears substantively only in US NIH-funded resources such as TCGA, CPTAC-3, and PLCO, while widely used challenge benchmarks including CAMELYON, PANDA, and TUPAC16 provide no demographic fields at all~\citep{komen2024batch}. TCGA's dominance compounds this problem: it recruits through US academic centres and comprises predominantly non-Hispanic White patients~\citep{wang2018tcga}. The dataset is simultaneously the richest demographic source and the strongest site-leakage source in the corpus. This creates a structural risk of confounding demographic disparity with site--demographic correlation~\citep{howard2021site, vaidya2024demographic, celi2022sources}. The geographic concentration is pronounced: $\sim$84\% (66/78) of first/corresponding authors are from North America and Europe, with East Asia contributing most of the remainder. The corpus focuses on cancers prevalent in high-income countries, while cancers concentrated in low-income countries are underrepresented. Cervical cancer appears in only 2 studies and paediatric cancer in none, making these populations virtually absent from the fairness literature~\citep{celi2022sources, musa2026cross, ceccon2025underrepresentation}. Models deployed beyond their training geography show performance degradation of 10--25\% across populations~\citep{musa2026cross}, and cross-population mutation-rate variation further amplifies diagnostic disparities in underrepresented settings~\citep{polin2025mutation}.

Federated learning offers a partial remedy. Prop-FFL reduced inter-hospital variance from 32.54 to 10.00~\citep{hosseini2023proportionally}, and federated learning-trained melanoma classifiers outperformed centralised baselines (AUROC 0.913 versus 0.905)~\citep{haggenmuller2024federated}. But federated learning cannot manufacture diversity the federation does not possess~\citep{celi2022sources}, and fairness-aware aggregation still requires client-side demographic metadata that is typically absent.

\subsubsection{Benchmark and Reporting Gaps}
\label{sec:discussion:gaps}
Three infrastructure deficits prevent fairness results from accumulating across studies. First, no histopathology-specific fairness benchmark exists that is simultaneously WSI-native, multi-site with preserved-site cross-validation, annotated with demographic metadata, and spans both high- and low-prevalence cancers. MedFair contains no WSI cohort~\citep{zong2023medfair}, and WILDS-CAMELYON17 carries no demographic metadata~\citep{koh2021wilds, rui2025ethics}. Second, fairness metric reporting remains severely underdeveloped: only 19.2\% of studies report any fairness metric, 7.7\% report worst-group accuracy, two studies report expected calibration error, and none report subgroup calibration gap or attention-level fairness probes for MIL pipelines. We propose a six-item minimum reporting floor as a precondition for claims of bias mitigation: aggregate AUROC on in-distribution and out-of-distribution splits, a subgroup disparity metric, worst-group accuracy or AUROC, expected calibration error and subgroup calibration gap, a preserved-site leakage probe, and an attention-level fairness summary for MIL models. Third, no experimental paper in the corpus cites compliance with CLAIM or TRIPOD+AI, and three histopathology-specific extensions are needed covering tissue provenance, pre-processing audit, and fairness evaluation~\citep{sounderajah2021stardai, howard2021site, komen2024batch}. These deficits are infrastructural rather than epistemic. Cross-domain evidence confirms their scale: a scoping review of 692 FDA-approved AI/ML medical devices from 1995--2023 found that only 3.6\% reported race or ethnicity and 99.1\% provided no socioeconomic data, establishing that pre-market authorisation under existing frameworks has not enforced demographic transparency~\citep{muralidharan2024fda}. The exhortation to report fairly has been voiced for half a decade without reshaping practice; benchmarks, editorial policies, and reviewing norms are the binding instruments required to change behaviour. 

\subsubsection{Regulatory Compliance Gaps}
\label{sec:discussion:regulation}
Regulatory frameworks now mandate subgroup performance evaluation for high-risk medical AI. The FDA's Good Machine Learning Practice principles require representative training data and performance evaluation across relevant subgroups, with enforcement ramping up through 2025--2027 and non-compliance potentially resulting in Complete Response Letters delaying approval by 6--18 months. The EU AI Act (2024/1689) classifies pathology AI as high-risk, requiring fundamental rights impact assessments and ongoing post-market monitoring, with penalties up to €35M or 7\% of global annual turnover for non-compliance~\citep{euaiact2024}. The UK MHRA similarly requires fairness and bias mitigation as part of the total product lifecycle with mandatory incident reporting. However, the histopathology corpus lacks the infrastructure for compliance: no study reports pre-specified longitudinal fairness-drift monitoring, demographic metadata is absent from most datasets, and uncorrected site leakage persists. While precedents exist---the FDA required extensive subgroup analysis for IDx-DR approval and delayed Google Health's mammography AI pending additional demographic data---no histopathology AI has faced regulatory action for fairness violations as of March 2026. Regulatory heterogeneity creates arbitrage risk, necessitating harmonized international standards through bodies like the International Medical Device Regulators Forum, whose AI/ML guidance expected in 2026--2027 should align enforcement mechanisms across jurisdictions.

\subsection{Study Limitations}
\label{sec:discussion:limitations}
This review's histopathology-specific scope permits technical granularity (TSS AUROC regimes, linear-probe recovery rates, WILDS-CAMELYON17 gaps) that broader medical AI fairness reviews must abstract away. Its analytical base of 78 experimental studies with paper-specific findings quantifies the reporting--practice gap rather than merely asserting it. Structured, auditable extraction via JSON schemas and a 39-dataset metadata audit enables any claim to be traced to primary data~\citep{sounderajah2021stardai}, and the integrated analytic framework combining an eight-category bias taxonomy, sixteen-family mitigation grading, and dataset-provenance audit provides a reusable foundation for future systematic updates~\citep{chen2023algorithmic, montezuma2025unbiased}. These strengths carry trade-offs. An English-language restriction may exclude relevant work from non-Anglophone jurisdictions, compounding the Global South under-representation this review documents~\citep{musa2026cross, celi2022sources}. We identified 23 potentially relevant non-English studies (17 Chinese, 4 Spanish, 2 German) during screening; abstract review suggested most focused on technical image processing without fairness analysis, but we cannot exclude the possibility of relevant unpublished work. Inclusion of preprints, necessary to capture the foundation-model literature, introduces the risk that some reported effect sizes may not survive peer review~\citep{komen2024batch, huang2025knowledge}. Single-reviewer extraction with 20\% audit means residual misclassification at bias-category boundaries cannot be excluded (inter-rater agreement: 88.9\%, Cohen's $\kappa=0.85$), and qualitative effectiveness grading rather than formal meta-analysis, while justified by the heterogeneity of study designs, limits quantitative pooling.

Two additional limitations warrant emphasis. First, \textbf{publication bias} may inflate reported fairness gaps and mitigation effectiveness. Studies finding no significant fairness disparities, or reporting failed mitigation attempts, may be less likely to submit or pass peer review. The 78-study corpus likely over-represents positive findings (bias detected, mitigation successful) relative to null results. We attempted to mitigate this by including grey literature and preprints, but cannot exclude the possibility that the true prevalence of fairness metrics in the field is higher than our 19.2\% estimate if studies without fairness findings remain unpublished.

Second, the \textbf{multiple comparisons problem} affects our bias category assignments. With 78 studies and 8 bias categories (plus multiple mitigation families and outcome metrics), some associations may arise by chance. We mitigated this risk by requiring explicit experimental evidence for each bias category assignment, using conservative evidence grading (strong/moderate/mixed/failed), and focusing on effect sizes rather than statistical significance. Nevertheless, readers should interpret individual category counts (e.g., $n=5$ for demographic bias, $n=7$ for institutional bias) as descriptive rather than inferential statistics. The 95\% confidence intervals reported throughout the results section quantify sampling uncertainty but do not adjust for multiple testing across bias categories.

Most importantly, the review inherits the same distributional distortions it critiques: TCGA dominance, high-income geographic concentration, and structural absence of demographic metadata. Every corpus-wide statistic should be read as valid within the sampling frame of evidence available to March~2026 rather than as a global estimate, and the search cut-off excludes second-generation foundation models and post-cut-off regulatory guidance from this synthesis~\citep{biswas2026digital}. Structural claims (site leakage unresolved, fairness under-reported, TCGA over-anchoring, infrastructure binding) are supported by convergent evidence across multiple dataset families. Quantitative claims rest on a smaller evidence base and should be cited with their empirical provenance.

\subsection{Recommendations and Research Priorities}
\label{sec:discussion:agenda}

\subsubsection{Reporting Checklist and Benchmark Proposal}
\label{sec:discussion:agenda:checklist}
We propose two concrete responses to the gaps identified above. The first is a 24-item reporting checklist organised into five blocks (Table~\ref{tab:checklist}), designed to complement CLAIM, TRIPOD+AI, and STARD-AI~\citep{sounderajah2021stardai}. The second is a Fairness Benchmark for Histopathology (FB-Histo) built on six pillars: multi-site federation (10+ sites across 4 continents), demographic stratification with mandatory patient-level annotation, a task suite covering cancer-subtype classification through biomarker prediction, canonical metrics aligned with Block D, preserved-site cross-validation as default, and governance with pathologist, patient, and Global South representation~\citep{howard2021site, vaidya2024demographic, zong2023medfair}. Governance would follow a multi-stakeholder consortium model with 9-member steering committee (including 2 Global South representatives with veto power), Fairness Review Board, and Technical Committee. Tiered data access (open, controlled, federated) balances openness with privacy, while incentive mechanisms including co-authorship for data contributors, cash prizes for competition participants, and subsidised compute credits for Global South participation encourage broad engagement. This model builds on established pathology benchmarks while addressing their demographic representativeness limitations.

\begin{table}[htbp]
  \centering
  \caption{Proposed 24-item reporting checklist for fairness in histopathology AI,
           organised into five blocks. Block~D operationalises the six-item
           minimum metric floor proposed in this review.}
  \label{tab:checklist}
  \small
  \begin{tabularx}{\linewidth}{@{}l l X l@{}}
    \toprule
    \textbf{Block} & \textbf{Focus} & \textbf{Required items} & \textbf{Ref.} \\
    \midrule
    A & Tissue provenance
      & Sites, scanners, magnification, staining protocol, quality-control log
      & \citep{howard2021site, komen2024batch} \\[2pt]
    B & Pre-processing audit
      & Normaliser identity, template slide, augmentation pipeline
      & \citep{lin2025stain, voon2023evaluating} \\[2pt]
    C & Cohort demographics
      & Per-subgroup sample sizes, attribute source (self-reported vs.\ inferred),
        geographic composition, missingness disclosure, consent provisions
      & \citep{vaidya2024demographic, celi2022sources} \\[2pt]
    D & Fairness evaluation
      & Subgroup AUROC (ID \& OOD), subgroup disparity metric, worst-group
        accuracy/AUROC, ECE and subgroup calibration gap, preserved-site
        leakage probe, MIL attention-level fairness summary
      & \citep{zong2023medfair, vaidya2024demographic, roschewitz2023automatic, howard2021site, chen2023algorithmic} \\[2pt]
    E & Reproducibility
      & Code availability, model weights, dataset access tiers,
        evaluation-harness documentation
      & \citep{koh2021wilds, zong2023medfair} \\
    \bottomrule
  \end{tabularx}
\end{table}

\subsubsection{Methodological Research Priorities}
Beyond these infrastructure deliverables, four methodological research priorities emerge. Fairness-aware foundation-model pretraining is the most pressing because current models encode TSS at over 90\% accuracy~\citep{komen2024batch}, and post-hoc correction addresses individual tasks without resolving the contaminated substrate. The goal is to incorporate batch- and demographic-invariance constraints into pretraining itself through site-adversarial self-supervision, scanner-conditioned disentanglement, or contrastive objectives that treat same-patient, different-scanner slides as positive pairs~\citep{huang2025knowledge, lin2025contrastive}. A second priority is causal and counterfactual modeling. Nearly all validated methods operate on associations and cannot distinguish spurious TSS signatures from biological signal~\citep{howard2021site, komen2024batch}. Cross-modal evidence confirms this ceiling: across radiology, MRI, and fundus imaging, pre-training improves accuracy but has minimal impact on fairness (p$>$0.05 in 18 of 24 settings), with no method consistently outperforming empirical risk minimisation across demographic groups~\citep{pham2025boundaries}. Within histopathology specifically, \citet{correamedero2024causal} demonstrated equalised survival performance across seven sites using causal adjustment, while \citet{shabazian2025joint} coupled fairness with explanation-guided training.

The third and fourth priorities concern implementation mechanisms. Federated learning requires fairness-aware aggregation beyond Prop-FFL~\citep{hosseini2023proportionally}, encompassing federated fairness auditing without centralising data and privacy-preserving substrates calibrated to avoid erasing subgroup disparities~\citep{lu2022federated}. FlexFair addresses this challenge directly by jointly optimising demographic parity, equalised odds, and accuracy without centralising data~\citep{xing2025flexfair}, and should be integrated with fairness-aware pretraining objectives. Generative models offer a practical route to augmenting underrepresented subgroups: \citet{ktena2024generative} achieved 48.5\% relative performance improvement via conditional diffusion (absolute fairness improvement of $\pm$30.0\%), and \citet{vanbooven2025mitigating} improved Gleason grading accuracy via synthetic data augmentation~\citep{ktena2024generative, vanbooven2025mitigating}. These four priorities are mutually reinforcing: fairness-aware pretraining provides the substrate that causal interventions require, federated learning supplies the multi-institutional diversity both depend on, and generative augmentation fills the long-tail intersectional cells that neither training data nor federation alone can supply.

\subsubsection{Policy, Implementation, and Ethical Considerations}
\label{sec:discussion:policy}
Funding bodies should condition grants on representativeness plans with demographic and geographic targets, earmarking budgets for Global South sites~\citep{celi2022sources, biswas2026digital}. Journals should require checklist completion (Blocks A--E) as a submission precondition~\citep{sounderajah2021stardai, chen2023algorithmic}, while challenge organisers should extend benchmarks like MIDOG and GrandQC to incorporate fairness evaluation~\citep{aubreville2022mitosis, weng2024grandqc}. Demographic metadata collection must balance audit necessity against privacy risk: consent should be purpose-specific, collection limited to minimum necessary attributes, and re-identification authorised only for adverse-event investigation~\citep{norori2021addressing, chen2023algorithmic}. Self-reported race, genetic ancestry, and tissue-source site address different questions and require joint deliberation among clinicians, patients, and ethicists~\citep{howard2021site, vaidya2024demographic}. Hybrid-responsibility frameworks are needed for foundation-model failures, as pathologists cannot detect embedding-level site leakage through slide inspection alone~\citep{cheng2021challenges, komen2024batch}. Finally, equitable care requires international consortia and capacity-building partnerships to ensure populations absent from training data are not excluded from benefits~\citep{musa2026cross, norori2021addressing}.
